\section*{Extended Synopsis}

The primary objective of the \textbf{QUPROB} project is to establish a unified theoretical and experimental framework for controlling non-equilibrium phase transitions in topological quantum systems under stochastic environmental influences. By bridging non-Markovian open quantum systems theory with topological shielding, QUPROB aims to demonstrate a new class of dynamically protected logical gates that remain robust under realistic, non-white environmental noise.

\pbreak
State-of-the-art quantum computing architectures rely on passive topological protection, which assumes static error models or isolated systems. However, real-world devices operate under dynamic, non-Markovian (memory-bearing) colored noise, which rapidly degrades topological order. QUPROB resolves this limitation through three key pillars:
\begin{enumerate}[label=\textbf{\arabic*.}, leftmargin=2em]
    \item \textbf{Theoretical Foundation:} Developing a time-dependent stochastic field theory for open topological systems.
    \item \textbf{Dynamical Control:} Leveraging stochastic resonance to actively stabilize topological edge states.
    \item \textbf{Experimental Validation:} Implementing these protocols on a 10-qubit superconducting transmon platform.
\end{enumerate}

\pbreak
\begin{highlightbox}
\textbf{Key Project Output:} The successful execution of QUPROB will result in a $10\times$ improvement in coherence times for topological logical states under strong colored noise, laying the groundwork for fault-tolerant quantum computing without prohibitive error-correction overheads.
\end{highlightbox}

\vspace{1.0em}
